% Based on arXiv:2404.19737, "Better & Faster Large Language Models via Multi-token Prediction", CC BY 4.0.
% https://arxiv.org/abs/2404.19737
% The prose is from the paper, with two edits: the closing `$` after "n=6" is
% removed, leaving a stray `$` partway through the document, and a repeated
% "the" is added after it. The repeated word must still be linted.
\documentclass{article}

\begin{document}

\subsection{Searching for the optimal $n$}
\label{sect:optimal-future}

\begin{figure}[t!]
    \centering
    \includegraphics[width=\linewidth]{img/backward_order.png}
    \caption{\textbf{Order of the forward/backward in an $n$-token prediction model with $n = 2$ heads.} By performing the forward/backward on the heads in sequential order, we avoid materializing all unembedding layer gradients in memory simultaneously and reduce peak GPU memory usage.}
    \label{fig:backward}
\end{figure}

To better understand the effect of the number of predicted tokens, we did comprehensive ablations on models of scale 7B trained on 200B tokens of code. Results in table \ref{tab:future_tokens_results} show that training with 4-future tokens outperforms all the other models consistently throughout HumanEval and MBPP for pass at 1, 10 and 100 metrics. Interestingly, for APPS/Intro, $n=6 takes the lead with +0.7\%, +3.0\% and +5.3\%.

It is very likely that the the optimal window size depends on input data distribution. As for the byte level models the optimal window size is more consistent (8 bytes) across these benchmarks.

\end{document}
